\section{Ampliación de la sesión de Q\&A: limitaciones que las slides de la clase no mencionan}

La sesión de Q\&A no añadió ningún visual state independiente, pero completó varias acotaciones que determinan si las conclusiones técnicas son confiables: el prefill cost del long context, la relación de suma no nula entre data y architecture, la incertidumbre sobre la ventaja de diffusion, el origen de la alta resolución de CogAgent y si un modelo de video puede aprender física a partir de la raw signal.

\subsection{Long context, data y architecture: tres equilibrios}

Primero, un contexto largo puede reducir la necesidad de un retrieval pipeline externo, pero la prefill latency y la interferencia de tokens irrelevantes siguen presentes. Segundo, los data pueden expresar muchos task-specific biases, pero una general architecture improvement todavía puede cambiar la compute frontier de todas las tareas. Tercero, colocar todos los documents o reasoning branches en el context exige que el modelo aprenda a aprovecharlos; no basta con ampliar la ventana.

\begin{importantbox}{Conclusiones de ingeniería de la sesión de Q\&A}
\begin{enumerate}
  \item Para documentos largos con respuestas cortas, puede ser razonable aceptar un prefill más lento; para un agent en tiempo real, no necesariamente.
  \item Ante un comportamiento específico, conviene intentar primero con data/evaluation; ante un cuello de botella que se repite entre tareas, considerar después la architecture.
  \item Si el reasoning procedure incluye rutas fallidas, colocarlas explícitamente en el context puede aportar más información que conservar solo el beam score, aunque esto todavía requiere experimentos.
\end{enumerate}
\end{importantbox}

\teachervoice{01:18:10--01:19:59, el ponente se inclina por colocar las alternative paths junto con la failure information en el context, para que el modelo aprenda a compararlas, en lugar de usar solo una hard-coded beam probability; es una observación empírica, no un universal theorem.}

\subsection{La video understanding no adquiere automáticamente el mundo físico}

Los VLM/video models actuales dependen sobre todo de text-image o text-video pairs. Si la training signal proviene principalmente de annotation humana, el modelo aprende los fenómenos que cubren las descripciones humanas, y no necesariamente recupera a partir del raw motion la conservation, el contact, la object permanence o las causal interventions.

\begin{warningbox}{«Haber visto muchos videos» no equivale a «haber aprendido física»}
Para aprender física a partir de video sin anotar pueden hacer falta temporal prediction, masked video modeling, world-model objectives, interaction data o una simulation verificable. Ampliar solo los annotated pairs todavía puede llevar a aprender shortcuts del lenguaje y correlaciones superficiales.
\end{warningbox}

\teachervoice{01:16:15--01:18:10, el ponente afirma con claridad que, si la fuente de datos no contiene physical rules, el preentrenamiento actual basado en pairs no las aprenderá de forma automática; considera que la nueva video self-supervision es un problema de investigación necesario.}

\subsection{Resumen de la sección}

La sesión de Q\&A convierte los headlines de la exposición principal en juicios de ingeniería con condiciones: el context tiene latencia, los data no eliminan la architecture, diffusion no ha superado matemáticamente a autoregression, una GUI demo no representa un agent confiable y la video quantity tampoco equivale a physical understanding.
